\documentclass{sprawozdanie-agh}

\usepackage[utf8]{inputenc}
\usepackage{matlab-prettifier}
\usepackage{siunitx}
\usepackage{booktabs,caption,ragged2e}

\makeatletter

% ========= Początek Dokumentu ========= 
\begin{document}
% ========= Strona tytułowa ========= 
\przedmiot{Teoria Sterowania II: Ćwiczenia Laboratoryjne}
\tytul{Wytyczne dla Ćwiczenia 5}
\podtytul{Optymalizacja Parametryczna}
\kierunek{Automatyke i~Robotyka}
\autor{Dariusz Cieślar, Maciej Różewicz}
\data{Kraków, 3 marca 2026}

\stronatytulowa{}
% ========= Treść Właściwa =========
\section{Cel ćwiczenia}
Celem ćwiczenia jest zapoznanie się z~optymalizacją parametryczną układów sterowania.
Dotyczy ona problemów, gdzie struktura regulatora i~model matematyczny obiektu są stałe, a~optymalizowane są jedynie wartości niektórych parametrów.
Optymalność tę można rozumieć na różne sposoby, a~dobór odpowiedniej funkcji celu, która jest optymalizowana, jest zagadnieniem zasadniczym.
Zwykle żąda się, aby układ regulacji możliwie dokładnie i~szybko odtwarzał wartość zadaną, z~drugiej zaś strony, by był możliwie odporny na sygnały zakłócające.
Często wymagania te są sprzeczne i~poszukiwać trzeba pomiędzy nimi kompromisu.

%%%%%%%%%%%%%%%%%%%%%%%%%%%%%%%%%%%%%%%%%%%%%%%%%%%%%%%%%%%
\section{Przebieg ćwiczenia}
Dla każdego z~podanych systemów w~trakcie analizy należy dla układów 1 i~2:
\begin{itemize}
    \item utworzyć model symulacyjny układu sterowania w~Simulinku (łącznie z~wejściami zakłóceń),
    \item wyznaczyć analitycznie ograniczenia na zakres zmiennych projektowych (w zależności od zadanej struktury regulatora) - zaznaczyć ten obszar na płaszczyźnie,
    \item wykonać numerycznie optymalizacje ze względu na wskaźnik jakości $J_{1}$,
    \item wykonać numerycznie optymalizację ze względu na wskaźnik jakości $J_{7}$,
    \item wykonać numerycznie optymalizację ze względu na odporność na obciążenie,
    \item wykonać numerycznie optymalizację ze względu na wskaźniki $J_{1}$  i~na obciążenie,
    \item dla każdego zestawu dobranych nastaw regulatora wykonać symulacyjnie:
        \begin{itemize}
            \item odpowiedź na skok jednostkowy bez zakłóceń (pokazać wartości sterowania),
            \item odpowiedź na skok jednostkowy z~zakłóceniem obciążeniowym (skok jednostkowy, ale w~innej chwili niż wartość zadana)(pokazać wartości sterowania).
        \end{itemize}
\end{itemize}

    %======================================================
    \subsection{Układ 1}
    Dla obiektu opisywanego transmitancją:
    \begin{equation}
        \label{eq:system_1}
        G_{O}(s) = \frac{2s+3}{s^3+2s^2+s},
    \end{equation}
    w~układzie regulacji automatycznej, rozważyć do optymalizacji parametrycznej regulatory:
    \begin{itemize}
        \item P: $G_{R,P}(s) = K_{P}$,
        \item PD: $G_{R,PD}(s) = K_{P}(1 + T_{D}s)$.
    \end{itemize}

    %======================================================
    \subsection{Układ 2}
    Dla obiektu opisywanego transmitancją:
    \begin{equation}
        \label{eq:system_2}
        G_{1}(s) = \frac{1}{s^2+0.5s+1},
    \end{equation}
    w~układzie regulacji automatycznej, rozważyć do optymalizacji parametrycznej regulatory:
    \begin{itemize}
        \item P: $G_{R,P}(s) = K_{P}$,
        \item PI: $G_{R,PD}(s) = K_{P}(1 + \frac{1}{T_{D}s})$.
    \end{itemize}

    %======================================================
    \subsection{Układ 3}
    Wyznaczyć optymalny regulator proporcjonalny $G_{R}(s) = K$, dla obiektu opisywanego transmitancją:
    $$
        G_{O}(s) = \frac{1}{s^2 + s + 1},
    $$
    minimalizujący czas narastania, przy przeregulowaniu nie przekraczającym $10\%$.

\end{document}
